\documentclass{article}
\usepackage[T1]{fontenc}
\usepackage{lmodern}
\usepackage{graphicx}
\usepackage{amsmath,amssymb}
\usepackage[a4paper,margin=2cm]{geometry}
\usepackage[absolute,overlay]{textpos}
\setlength{\TPHorizModule}{1mm}
\setlength{\TPVertModule}{1mm}
\usepackage[indonesian]{babel}
\usepackage{hyperref}
\setlength{\emergencystretch}{2em}
% unit: Halaman 93

\begin{document}

% === Halaman 93 (llm) ===

\begin{itemize}
    \item Salah satu cara memperbesar faktor kerja untuk \textit{exhaustive key search}: adalah dengan cara enkripsi tidak dilakukan terhadap huruf individual, tetapi dalam blok huruf.
    \item Misal, pesan \texttt{kriptografi} dipecah menjadi kelompok 4-huruf:
    \[ \texttt{krip togr afi} \]
\end{itemize}

(ekivalen dengan 10170815 19140617 000508, dengan memisalkan 'A' = 0, 'B' = 1, ..., 'Z' = 25)

\end{document}
